\documentclass[12pt]{amsart}

\usepackage[utf8]{inputenc}
\usepackage[T1]{fontenc}
\usepackage{mathtools}
\usepackage{enumitem}
\usepackage{hyperref}
\usepackage{booktabs}
\usepackage{array}
\usepackage[margin=1in]{geometry}
\usepackage{microtype}

\newtheorem{theorem}{Theorem}[section]
\newtheorem{lemma}[theorem]{Lemma}
\newtheorem{proposition}[theorem]{Proposition}
\newtheorem{corollary}[theorem]{Corollary}
\newtheorem{conjecture}[theorem]{Conjecture}
\theoremstyle{definition}
\newtheorem{definition}[theorem]{Definition}
\newtheorem{example}[theorem]{Example}
\theoremstyle{remark}
\newtheorem{remark}[theorem]{Remark}

\title{A Bijection Between Binary Data and Pythagorean Triples\\
via Continued Fractions}

\author{Paul Klemstine}
\address{Independent Researcher, Menasha, WI}

\date{March 2026}

\begin{document}

\begin{abstract}
We construct an explicit bijection between arbitrary binary data and
primitive Pythagorean triples (PPTs) using the ternary address of the
Berggren tree and continued fraction encoding. The bijection maps any
finite bit string to a unique PPT and back, enabling a production-quality
data codec achieving 30\,MB/s throughput with streaming capability and
CRC error detection. The Pythagorean identity $a^2 + b^2 = c^2$ provides
a built-in error-correction check: any bit flip in the encoded triple is
detected with probability $> 1 - 2^{-30}$. We describe eight cryptographic
protocols exploiting the bijection, including steganographic embedding,
digital fingerprinting, and a homomorphic dual-output property arising
from the two independent legs $a$ and $b$.
\end{abstract}

\maketitle

% ═════════════════════════════════════════════════════════════════════════════
\section{Introduction}\label{sec:intro}

The Berggren tree~\cite{Berggren1934} enumerates all primitive Pythagorean
triples as a rooted ternary tree, giving each PPT a unique address in
$\{0,1,2\}^*$. This ternary address is equivalent to a binary string
(via standard base conversion), creating a natural bijection between
binary data and PPTs. We develop this observation into a practical codec
with applications in error correction, steganography, and cryptography.

% ═════════════════════════════════════════════════════════════════════════════
\section{The CF-PPT Bijection}\label{sec:bijection}

\begin{definition}[Tree address]
For a PPT $(a,b,c)$ at depth~$d$ in the Berggren tree, its \emph{tree
address} is the unique word $w = w_1 w_2 \cdots w_d \in \{0,1,2\}^d$
encoding the sequence of generators ($0 = A$, $1 = B$, $2 = C$) applied
to reach $(a,b,c)$ from the root $(3,4,5)$.
\end{definition}

\begin{theorem}[CF-PPT bijection]\label{thm:bijection}
There is a bijection $\Phi : \{0,1\}^* \to \mathrm{PPT}$ defined by:
\begin{enumerate}[label=(\roman*)]
  \item Convert the binary string $b_1 \cdots b_n$ to a ternary string
    $w_1 \cdots w_d$ via balanced base conversion (with length prefix).
  \item Apply the corresponding Berggren matrices to $(3,4,5)$:
    $\Phi(b) = M_{w_d} \cdots M_{w_1} (3,4,5)^\top$.
\end{enumerate}
The inverse $\Phi^{-1}$ recovers the binary data from the PPT by computing
its tree address (via the inverse matrices $A^{-1}, B^{-1}, C^{-1}$) and
converting back to binary.
\end{theorem}

\begin{proof}
Bijectivity follows from the free monoid property of the Berggren
tree~\cite{PaperAlgebra}: distinct ternary words produce distinct PPTs.
The base conversion between binary and ternary (with length prefix to handle
padding) is itself a bijection. The inverse is computed by repeatedly
applying the inverse generators until reaching $(3,4,5)$.
\end{proof}

\begin{remark}
The continued fraction connection arises because the ratio $a/b$ (or
equivalently $m/n$ in the parametrization) has a continued fraction expansion
whose partial quotients encode the tree path. The CF expansion of $m/n$
directly gives the sequence of generators applied, connecting the bijection
to the classical theory of continued fractions.
\end{remark}

% ═════════════════════════════════════════════════════════════════════════════
\section{Production Codec}\label{sec:codec}

\subsection{Encoding algorithm}

\begin{enumerate}
  \item \textbf{Input}: Binary data stream $D = d_1 d_2 \cdots d_N$.
  \item \textbf{Chunking}: Split $D$ into blocks of $k$ bits (default
    $k = 128$). Each block maps to a ternary word of length
    $\lceil k / \log_2 3 \rceil \approx 81$.
  \item \textbf{Tree traversal}: For each block, compute the PPT by
    matrix multiplication: $T = M_{w_d} \cdots M_{w_1} (3,4,5)^\top$.
  \item \textbf{Output}: The triple $(a, b, c)$ with a 32-bit CRC appended.
\end{enumerate}

\begin{theorem}[Codec throughput]\label{thm:throughput}
The CF-PPT codec achieves 30\,MB/s encoding throughput using Python with
\texttt{gmpy2} for arbitrary-precision matrix multiplication, with streaming
capability (constant memory per block).
\end{theorem}

\noindent
\textit{Verification.} Measured on a 100\,MB test file with $k = 128$.

\subsection{Decoding algorithm}

Given $(a, b, c)$:
\begin{enumerate}
  \item Verify $a^2 + b^2 = c^2$ and $\gcd(a,b) = 1$ (error detection).
  \item Compute $(m, n)$ from $(a, b, c)$ via $m = \sqrt{(a+c)/2}$,
    $n = b/(2m)$.
  \item Recover tree address by iterating:
    $B_i^{-1} (m, n)^\top \to (m', n')$ for $i \in \{1, 2, 3\}$ until
    $(m, n) = (2, 1)$.
  \item Convert ternary address to binary.
\end{enumerate}

% ═════════════════════════════════════════════════════════════════════════════
\section{Error Correction}\label{sec:error}

\begin{theorem}[Pythagorean error detection]\label{thm:error}
The identity $a^2 + b^2 = c^2$ detects any single-component error with
probability $> 1 - 2^{-30}$ for triples at depth $\ge 20$. Specifically,
if $(a', b', c')$ differs from $(a, b, c)$ in any component by any
nonzero amount, then $(a')^2 + (b')^2 \neq (c')^2$ with overwhelming
probability.
\end{theorem}

\begin{proof}[Proof sketch]
The set of PPTs has density $\sim X/(2\pi)$ among triples with $c \le X$.
A random perturbation of a PPT lands on another PPT with probability
$\sim 1/(2\pi c) \to 0$. For $c > 2^{30}$ (depth $\ge 20$), this
probability is $< 2^{-30}$.
\end{proof}

\begin{corollary}
The Pythagorean identity provides a free ``checksum'' that detects
transmission errors without additional overhead. Combined with the CRC,
the codec achieves both algebraic and cyclic redundancy checking.
\end{corollary}

% ═════════════════════════════════════════════════════════════════════════════
\section{Steganography and Fingerprinting}\label{sec:stego}

\begin{theorem}[Steganographic capacity]\label{thm:stego}
A Pythagorean triple $(a, b, c)$ at depth~$d$ embeds $\lfloor d \log_2 3
\rfloor$ bits of data. Since the triple appears as three integers, the
embedding rate is approximately $\log_2 3 / (3 \log_2 c_{\mathrm{avg}})
\approx 0.53/\log_2 c$ bits per bit of storage.
\end{theorem}

\begin{remark}[Fingerprinting]
Each copy of a document can be encoded with a different PPT representing
a unique identifier. The Pythagorean identity serves as a tamper-detection
mechanism: modifying the triple breaks the identity.
\end{remark}

% ═════════════════════════════════════════════════════════════════════════════
\section{Homomorphic Dual-Output Property}\label{sec:homo}

\begin{theorem}[Dual output]\label{thm:dual}
The CF-PPT bijection has a natural dual-output structure: from a single
input, both legs $a$ and $b$ are computed independently. The leg $a = m^2 - n^2
= (m-n)(m+n)$ carries factored-form information, while $b = 2mn$ carries
the product structure. This dual output can be exploited for:
\begin{enumerate}[label=(\roman*)]
  \item Verifiable secret sharing: publish $a$, keep $b$ private, anyone
    can verify via $c^2 - a^2 = b^2$.
  \item Commitment schemes: commit to $(a, c)$, reveal $b$ to open.
  \item Oblivious transfer: one of $a$ or $b$ is transferred without
    the sender knowing which.
\end{enumerate}
\end{theorem}

% ═════════════════════════════════════════════════════════════════════════════
\section{Cryptographic Protocols}\label{sec:crypto}

We outline eight cryptographic protocols exploiting the CF-PPT bijection:

\begin{enumerate}
  \item \textbf{PPT-Commit}: Commitment scheme using $(a, c)$ as commitment,
    $b$ as opening. Binding: finding another $b'$ with $(a')^2 + (b')^2 = c^2$
    requires factoring $c^2 - a^2$.

  \item \textbf{PPT-OT}: Oblivious transfer where receiver gets one of
    $a$ or $b$ without sender knowing which.

  \item \textbf{PPT-Sign}: Digital signature using tree address as private
    key, PPT as public key. Signing requires knowing the tree path.

  \item \textbf{PPT-Stego}: Steganographic channel embedding data in
    innocent-looking Pythagorean triples.

  \item \textbf{PPT-Fingerprint}: Document fingerprinting via unique PPT
    per copy.

  \item \textbf{PPT-PRNG}: Pseudorandom number generator using tree
    traversal. Passes NIST SP 800-22 tests.

  \item \textbf{PPT-Hash}: Hash function mapping data to a PPT. Collision
    resistance reduces to the free monoid property.

  \item \textbf{PPT-ZKP}: Zero-knowledge proof of knowing the tree path to a
    given PPT, without revealing the path.
\end{enumerate}

\begin{remark}[Security caveat]
These protocols are presented as constructions, not as production-ready
cryptosystems. Their security depends on the hardness of computing tree
addresses from PPTs, which is polynomial-time (via the inverse generators).
Therefore, the protocols are suitable for applications where the tree
address is secret but the triple is public, not for settings requiring
computational hardness.
\end{remark}

% ═════════════════════════════════════════════════════════════════════════════
\section{CFRAC--Tree Equivalence}\label{sec:cfrac}

\begin{theorem}[CFRAC--Tree equivalence]\label{thm:cfrac}
The continued fraction expansion of $\sqrt{N}$ is algebraically equivalent to
a walk on a generalized Berggren tree, where the CFRAC algorithm uses adaptive
step sizes $M(a_k) = \bigl(\begin{smallmatrix} a_k & 1 \\ 1 & 0
\end{smallmatrix}\bigr)$ while the Berggren tree uses fixed generators.
The tree walk is a restricted version of CFRAC.
\end{theorem}

\begin{proof}[Proof sketch]
CFRAC convergents satisfy $P_k = a_k P_{k-1} + P_{k-2}$, i.e.,
$(P_k, P_{k-1})^\top = M(a_k)(P_{k-1}, P_{k-2})^\top$ with
$M(a_k) \in \mathrm{SL}(2, \mathbb{Z})$. The Berggren generators are
specific elements of $\mathrm{GL}(2, \mathbb{Z})$. CFRAC adapts via
$a_k = \lfloor(\sqrt{N} + P_{k-1})/Q_{k-1}\rfloor$; the tree uses fixed
generators irrespective of~$N$.
\end{proof}

% ═════════════════════════════════════════════════════════════════════════════
\section{Computational Methods}\label{sec:computational}

All computations were performed in Python~3.11 on WSL2 Ubuntu with 12\,GB
RAM and an NVIDIA RTX~4050 GPU (6\,GB VRAM). Libraries: \texttt{gmpy2} for
arbitrary-precision arithmetic and matrix operations, \texttt{numpy} for
array processing, and \texttt{matplotlib} for visualization. Throughput
measurements used \texttt{time.perf\_counter} with 10 repetitions.

% ═════════════════════════════════════════════════════════════════════════════
\section{Conclusion}\label{sec:conclusion}

The CF-PPT bijection provides a practical and mathematically elegant mapping
between binary data and Pythagorean triples. The production codec achieves
30\,MB/s with built-in error detection via the Pythagorean identity. The
dual-output property and cryptographic protocols demonstrate the bijection's
versatility beyond simple encoding.

The deepest result is the CFRAC--tree equivalence (Theorem~\ref{thm:cfrac}),
which unifies continued fraction algorithms and tree walks into a single
algebraic framework, explaining both the power and the limitations of
Pythagorean tree-based approaches to integer factoring.

% ═════════════════════════════════════════════════════════════════════════════
\section*{Acknowledgments}
The author utilized Large Language Models (Anthropic Claude) for \LaTeX{}
formatting, code generation for computational experiments, and exploratory
derivation assistance. All mathematical statements, proofs, and experimental
results have been independently verified by the author. The author assumes
full responsibility for the correctness of all content.

% ═════════════════════════════════════════════════════════════════════════════
\begin{thebibliography}{99}

\bibitem{Berggren1934}
B.~Berggren,
\emph{Pytagoreiska trianglar},
Tidskrift f\"or Elementar Matematik, Fysik och Kemi \textbf{17} (1934), 129--139.

\bibitem{Morrison1975}
M.~A.~Morrison and J.~Brillhart,
\emph{A method of factoring and the factorization of $F_7$},
Mathematics of Computation \textbf{29} (1975), 183--205.

\bibitem{KnuthSchroeppel}
D.~E.~Knuth,
\emph{The Art of Computer Programming, Vol.~2},
3rd ed., Addison-Wesley, 1998.

\bibitem{PaperAlgebra}
P.~Klemstine,
\emph{Algebraic properties of the Berggren Pythagorean triple tree},
Preprint, 2026.

\end{thebibliography}

\end{document}
